\section{Conclusion}
\label{sec:conclusion}

We have made the link structure of a corpus a first-class part of a decoder-only
language model's attention: documents are packed by graph traversal, a block-sparse
mask keeps them causally isolated except along explicit edges, and where document $A$
links to document $B$ the tokens of $A$ from the link onward may read $B$. On
Wikipedia hyperlinks and on code imports across several languages, this lowers the
likelihood of link-dependent completions relative to the same tokens scored without
the target, leaves single-document capability unchanged, and beats matched-compute
controls that grant more attention without the link gating. Varying only the fraction
of edges the mask may use shows the benefit on code scaling with realized link
density, saturating early in training, and being dominated by the density exposed at
inference.

The design's purpose is larger than these prerequisites. Because the grant opens
exactly where a reference closes, and because causality is never relaxed, the model is
taught by the plain pretraining objective that writing a reference is what makes a
document readable. That is the ingredient a flat concatenation baseline cannot have
and a retrieval-augmented system must be engineered to supply: the possibility of a
model that fetches from its corpus \emph{by generating a citation}, with no retriever,
no fusion module, and no post-training to teach the behaviour. The models in this
paper are too small to test that directly; designing the measurement, and scaling to
a regime where generation is coherent enough to make it, is the work this paper sets
up. Beyond that: sequence-parallel training so long-document modalities such as arXiv
can be co-packed at all, the corrected multi-source diversity experiment, and richer
edge types where the graph is the author's own (personal knowledge bases, code with
call graphs rather than imports).
